\chapter{Introducción}

Las series de tiempo son fundamentales en campos como la economía, la meteorología y la salud. La evolución de las redes neuronales ha permitido mejorar la precisión de los modelos de predicción, entre los que destacan las arquitecturas LSTM \cite{hochreiter1997lstm} y los transformadores \cite{vaswani2017attention}. Esta tesis estudia ambas familias de modelos y su adaptación a problemas de regresión con salidas simples y múltiples.

\section{Planteamiento del problema}

Aunque los transformadores han obtenido resultados relevantes en distintas tareas de aprendizaje profundo, su aplicación a la regresión de series de tiempo continúa siendo un área activa de investigación. Por ello, resulta necesario analizar su eficacia frente a modelos recurrentes tradicionales mediante experimentos comparables y métricas cuantitativas.

\section{Justificación}

El trabajo busca contribuir al conocimiento sobre la optimización de modelos de predicción y ofrecer evidencia experimental que permita seleccionar arquitecturas adecuadas para aplicaciones prácticas en distintos sectores.

\section{Objetivo general}

Desarrollar y evaluar modelos de redes neuronales basados en LSTM y transformadores adaptados para tareas de regresión en series de tiempo, explorando distintas configuraciones y arquitecturas para optimizar su desempeño en predicciones de salida simple y múltiple.

\section{Objetivos específicos}

\begin{enumerate}
  \item Investigar y analizar las arquitecturas recurrentes LSTM y las arquitecturas basadas en transformadores para la predicción de series de tiempo.
  \item Implementar y adaptar transformadores para tareas de regresión en series de tiempo.
  \item Diseñar un esquema de preprocesamiento y postprocesamiento de datos que mejore la calidad de las predicciones.
  \item Comparar el desempeño de los modelos LSTM y transformadores mediante RMSE, MAE, MAPE y $R^2$.
  \item Evaluar la capacidad de predicción en escenarios de salida simple y salida múltiple.
  \item Analizar los resultados obtenidos y proponer mejoras en los modelos.
\end{enumerate}

\section{Estructura de la tesis}

El capítulo 2 presenta los fundamentos teóricos; el capítulo 3 revisa trabajos previos; el capítulo 4 describe la metodología; el capítulo 5 reúne los resultados y su análisis; finalmente, el capítulo 6 expone las conclusiones, limitaciones y líneas de trabajo futuro.
